\begin{center}
\LARGE \textbf{Обходим Великий Российский Файрвол}\\
\large Часть 1 из 5\\
\normalsize\textbf{В России официально установлена цензура Интернета. Но контролировать Интернет не так просто, как это может показаться с первого взгляда.}
\end{center}
\begin{figure}[h]
\center{\includegraphics[width=0.25\linewidth]{I2P}}
\end{figure}
\textbf{I2P} --- анонимная оверлейная сеть, использующая принцип чесночной маршрутизации, исходные коды которой распространяются на условиях нескольких свободных лицензий. В отличии от Tor, который в первую очередь направлен на доступ к сайтам обычного интернета (хотя в нем и существуют скрытые сервисы, аналогичные ипсайтам в I2P, а в I2P можно получить доступ к внешнему Интернету, используя аутпрокси), основной целью I2P является доступ именно к скрытым ресурсам --- ипсайтам. Ипсайт от обычного вебсайта отличает только его нахождение в сети I2P.
\begin{important}
Никогда не используйте анонимные сети с настройками, которые позволяют использовать Интернет в одном браузере в обход самой сети. В этом случае будет достаточно вставить любое изображение из внешнего Интернета, чтобы получить ваш реальный IP-адрес.
\end{important}
\subsection*{Установка}
Для установки посетите \url{http://i2p2.de} или установите пакет с помощью пакетного менеджера вашего дистрибутива.
\subsection*{Использование}
\begin{important}
Не используйте аутпрокси для передачи конфиденциальных данных, владелец аутпрокси видит трафик нешифрованным.
\end{important}
После установки настройте свой браузер на использование HTTP-прокси 127.0.0.1:4444 и посетите страницу \url{http://127.0.0.1:7657}. Перед вами консоль маршрутизатора I2P --- место, из которого можно управлять всеми настройками I2P.\\
Для начала перейдите в меню <<Настройки I2P>> (\url{http://127.0.0.1:7657/config}) и установите ограничения скорости в соответствии со скоростью вашего интернета.\\
Остальные настройки можно оставить по умолчанию. Подождите некоторое время для полноценной интеграции с сетью, после чего вы сможете полноценно пользоваться сетью. Роутер желательно не выключать, так как при его перезапуске потребуется придется повторить этот процесс.\\
В I2P отсутствуют корневые DNS-сервера, копия адресной книги хранится на каждом роутере. На \url{http://rus.i2p} вы можете найти дополнительный список подписок, который можно добавить в susidns.
\subsection*{Некоторые ипсайты}
\begin{enumerate}
\item \url{http://forum.i2p} --- главный форум, есть рускоязычный раздел.
\item \url{http://rus.i2p} --- русская I2P-вики.
\item \url{http://pastethis.i2p} --- pastebin-подобный ресурс.
\item \url{http://flibusta.i2p} --- зеркало Флибусты, крупной библиотеки электронных книг.
\item \url{http://tracker.rus.i2p} --- русский торрент-трекер в I2P.
\end{enumerate}
\vfill
\scriptsize Создано на основе материалов из <<Настольной книги анонима>> --- \url{anonhandbook.org}, \url{anonhandbook.i2p}, \url{oxgzwiiypou6udlp.onion}
\normalsize
